\hspace{3mm}

\centerline{\bf \Huge Вторая тренировочная Олимпиада Эйлера}

\hspace{3mm}

 \begin{flushright}
     \scriptsize{«Что вершит судьбу Киви в ЭлМШ?\\ Некий незримый Расул Владимирович или Анджур Аркадьевич,\\ подобно длани господней, парящей над миром? \\По крайней мере, истинно то, \\что Киви не властен даже над своей волей»}
 \end{flushright}
 
\problem{1} Анджур Аркадьевич раздал ученикам 47 бабуфохов так, что каждая девочка получила на один бабуфох больше, чем каждый мальчик. Затем Анджур Аркадьевич раздал тем же детям 74 фрутеллы так, что каждый мальчик получил на одну фрутеллу больше, чем каждая девочка. Сколько всего было детей?

\hspace{3mm}

\problem{2}Существуют ли 11 пятизначных чисел, дающие разные остатки от деления на 11, любые два из которых получаются друг из друга перестановкой цифр?

\hspace{3mm}

\problem{3}В однокруговом шахматном турнире (каждый шахматист играл с каждым по одной партии) участвовало $N$ шахматистов. По окончании турнира оказалось, что каждый участник выиграл столько партий, сколько ничьих было зафиксировано в партиях, в которых этот участник не играл. Могло ли $N$ равняться 51?

\hspace{3mm}

\problem{4}В треугольнике $KLM$ точки $P$ и $Q$ – середины сторон $KM$ и $KL$ соответственно. На медиане LP выбрана точка $R$, не лежащая на $MQ$. Оказалось, что $RM = 2PQ$. Докажите, что $KR = LM$.

\hspace{3mm}

\problem{5} Ильзира готовилась к экзамену по физике и от скуки стала выписывать все возможные последовательности ${a_1, a_2, \dotsc a_{100}}$ содержащие все натуральные числа от 1 до 100. У Ильзиры есть три любимых натуральных числа не превосходящих 100 (разные). Какое наименьшее число последовательностей ей нужно выписать, чтобы точно встретить такую, в которой будет либо 2 либо 3 ее любимых числа подряд?